\chapter{Navigator Runtime API Contract}

\section{Normal Client Path}

The intended user-facing Navigator shape is data ingestion plus a simple update
call:
\begin{equation*}
\text{push timestamped data} \quad \rightarrow \quad \texttt{Navigator::update()}.
\end{equation*}

The exact names remain implementation choices, but the preferred direction is:
\begin{itemize}
    \item typed ingestion such as \texttt{push\_imu(...)} and
    \texttt{push\_gnss(...)}, or a constrained \texttt{push\_data(...)} family;
    \item internal fixed-capacity buffers owned by the configured Navigator or
    sensor graph;
    \item \texttt{Navigator::update()} as the ordinary client call that
    orchestrates all eligible work.
\end{itemize}

\section{Explicit Stages}

For tests, profiling, and advanced embedded scheduling, expose explicit stages:
\begin{itemize}
    \item \texttt{process\_strapdown\_integration()} for high-rate nominal IMU
    integration;
    \item \texttt{process\_covariance()} for medium-rate STM/covariance
    propagation and history buffering;
    \item \texttt{process\_measurements()} for low-rate asynchronous aiding.
\end{itemize}
For the first INS/GNSS implementation, the measurement stage must include
loosely coupled GNSS position and velocity updates with the configured antenna
lever arm from \cref{eq:v1_gnss_pos_pred,eq:v1_gnss_vel_pred}.

The temporary Phase 4 \texttt{process\_measurements()} path is not the final
INS Navigator API. It exists to preserve the current GNSS-only behavior while
the propagation seam is introduced.

\section{V1 Cadence Rules}

The first INS Navigator should keep the timing model explicit:
\begin{itemize}
    \item IMU increments are ingested and integrated at the IMU rate.
    \item Covariance prediction runs at the configured medium-rate prediction
    cadence. It uses the latest nominal state, the elapsed prediction interval,
    and the interval-average specific force from
    \cref{eq:v1_interval_average_specific_force}.
    \item Aiding measurements are processed when their timestamp and sensor
    policy make them eligible for the current update.
    \item V1 may store history for future delayed-measurement support, but it
    does not extrapolate latent measurements to the current time.
\end{itemize}

\section{History Buffer Semantics}

The v1 implementation should make history products inspectable even before
delayed-measurement replay is implemented:
\begin{itemize}
    \item IMU history stores timestamped corrected increments or the raw
    increments plus the bias state needed to reproduce the correction.
    \item Nominal-state history stores timestamped
    $(\hat{\pb}_{eb}^{e},\hat{\vb}_{eb}^{e},\hat{\quat}_{b}^{e},
    \hat{\bb}_g,\hat{\bb}_a)$ samples.
    \item Prediction history stores the covariance prediction node:
    timestamp, $\Ph_k$, $\mathbf{Q}_{d,k}$, and the interval endpoints.
\end{itemize}
These buffers are a runtime/debugging contract for v1 and a foundation for v2
measurement-time replay; they are not allowed to silently change the v1 update
math.

\section{Ownership Boundaries}

\begin{center}
\begin{tabular}{p{0.30\textwidth}p{0.58\textwidth}}
\toprule
Owner & Responsibility \\
\midrule
Navigator & public orchestration API, buffer scheduling, stage ordering \\
Mechanization policy & nominal quaternion/PVA integration from IMU data \\
Filter prediction policy & STM, process noise, covariance prediction/history \\
Update policy & aiding update lifecycle, injection, reset \\
Sensor/emulator/app support & data production, runtime validation, desktop simulation inputs \\
\bottomrule
\end{tabular}
\end{center}

\begin{implbox}
The product config should continue to expose explicit graph aliases. A likely
future shape is \texttt{Mechanization}, \texttt{Propagation}, \texttt{Filter},
\texttt{NavigatorUpdate}, and \texttt{Navigator}, where \texttt{Propagation}
coordinates mechanization and filter prediction without exposing planet/gravity
details directly to \texttt{Navigator}.
\end{implbox}

\section{Equation-to-Code Map}

\begin{center}
\begin{tabular}{p{0.28\textwidth}p{0.22\textwidth}p{0.38\textwidth}}
\toprule
Implementation piece & Equations & Notes \\
\midrule
IMU input correction & \cref{sec:imu_input_contract} & Biases are applied before coning/sculling. \\
Coning/sculling & \cref{eq:v1_coning_two_sample,eq:v1_sculling_two_sample} & Tests must lock sub-sample ordering. \\
Quaternion propagation & \cref{eq:v1_discrete_earth_rate_attitude_increment,eq:v1_discrete_attitude_delta_quat,eq:v1_discrete_quat_update,eq:v1_discrete_quat_normalization,eq:v1_quaternion_sign_sanity} & Implements the nominal attitude advance. \\
Velocity/position propagation & \cref{eq:v1_gravity_policy_contract,eq:v1_discrete_velocity_update,eq:v1_discrete_position_update} & Gravity comes from the configured gravity policy. \\
Specific-force input to $F_k$ & \cref{eq:v1_interval_average_specific_force} & Reuses the same sculling-corrected increment as mechanization. \\
Build $F_k$ and $G_k$ & \cref{eq:v1_fk_matrix,eq:v1_gk_matrix} & No hidden derivation is allowed in code. \\
STM and process noise & \cref{eq:v1_phi_first_order,eq:v1_qd_integral,eq:v1_qd_series,eq:v1_qd_first_order} & First-order is the v1 baseline and must be validated. \\
Covariance prediction & \cref{eq:v1_covariance_prediction} & Uses the same $\Ph_k$ and $\mathbf{Q}_{d,k}$ products buffered by prediction history. \\
EKF correction & \cref{eq:v1_measurement_residual,eq:v1_innovation_covariance,eq:v1_kalman_gain,eq:v1_error_state_correction,eq:v1_joseph_covariance_update} & Measurement models own $\hat{\mathbf{z}}$, $H$, and $R$. \\
Injection/reset & \cref{eq:v1_quat_injection,eq:v1_reset_covariance_update,eq:v1_attitude_reset_block,eq:v1_full_reset_jacobian} & Injection and covariance reset are one convention-locked unit. \\
GNSS aiding & \cref{eq:v1_gnss_pos_pred,eq:v1_gnss_pos_h,eq:v1_gnss_vel_pred,eq:v1_gnss_vel_h,eq:v1_gnss_vel_perturbation} & Includes antenna lever arm for position and velocity. \\
\bottomrule
\end{tabular}
\end{center}
